%%%%%%%%%%%%%%%%%%%%%%%%%%%%%%%%%%%%%%%%%
% Beamer Presentation
% LaTeX Template
% Version 1.0 (10/11/12)
%
% This template has been downloaded from:
% http://www.LaTeXTemplates.com
%
% License:
% CC BY-NC-SA 3.0 (http://creativecommons.org/licenses/by-nc-sa/3.0/)
%
%%%%%%%%%%%%%%%%%%%%%%%%%%%%%%%%%%%%%%%%%

%----------------------------------------------------------------------------------------
%	PACKAGES AND THEMES
%----------------------------------------------------------------------------------------

\documentclass[10pt]{beamer}              % only frames

\usepackage{tikz}
%\usepackage{pgf}
\usepackage{pgfpages}
%\setbeameroption{show notes on second screen=right}
%\setbeameroption{show notes}
\usetikzlibrary{shapes,arrows}
\mode<presentation> {
  \usetheme{Madrid}
  \usecolortheme{dolphin}
}

\usepackage{graphicx} % Allows including images
\graphicspath{{figs/}} % Directory in which figures are stored
\usepackage{booktabs} % Allows the use of \toprule, \midrule and \bottomrule in tables
\usepackage{animate}
\usepackage{multimedia}
\usepackage{caption}
\usepackage{subcaption}
\usepackage{tcolorbox}
\usepackage{appendixnumberbeamer}

%----------------------------------------------------------------------------------------
%	TITLE PAGE
%----------------------------------------------------------------------------------------

\title[Methods computational Biology]{{\huge Methods and applications in computational cell biology and biophysics}} % The short title appears at the bottom of every slide, the full title is only on the title page

\author[H.Martinez-Seara \& L. Cwiklik]{Lecturers: \textbf{Dr.~Hector~Martinez-Seara \& Dr.~hab.~Lukasz~Cwiklik} Supervisor:~\textbf{Prof.~Pavel~Jungwirth}} % Your name
\institute[] % Your institution as it will appear on the bottom of every slide, may be shorthand to save space
{
Charles University - Faculty XXX, Prague, Czech Republic \\ % Your institution for the title page
\medskip
\textit{hseara@uochb.cas.cz, lukasz.cwiklik@jh-inst.cas.cz} % Your email address
}
\date{\today}

\begin{document}

%-----------------------------------------------------------
\begin{frame}
  \titlepage % Print the title page as the first slide
  \begin{center}
    \vspace{-1cm}
    \includegraphics[width=0.3\textwidth]{logos/UOCHB_IOCB-logo.pdf}
  \end{center}
\end{frame}

%-----------------------------------------------------------
\begin{frame}{Program} 
\tableofcontents [hideallsubsections]
\end{frame}

%-----------------------------------------------------------
\section{Overview of computational strategies applied to cell and macromolecular imaging}
%-----------------------------------------------------------
\begin{frame}{Program} 
\tableofcontents [sectionstyle=show/shaded,subsectionstyle=show/show/hide] 
\end{frame}
%-----------------------------------------------------------

\subsection{Introduction to pattern and morphology recognition}

\subsection{Image manipulation methods}

\subsection{Concepts of probability theory applied to pattern recognition}

\subsection{Choosing an appropriate model for a given biological system and phenomenon}

\subsection{Decision theory as methodology behind pattern and morphology recognition techniques}

\subsection{Machine learning applied to image recognition}

\subsection{Data mining and high-throughput data analysis}

%-----------------------------------------------------------
\section{Overview of computational methods in  biological problems: from atomistic to macroscopic approaches}
%-----------------------------------------------------------
\begin{frame}{Program} 
\tableofcontents [sectionstyle=show/shaded,subsectionstyle=show/show/hide] 
\end{frame}

\subsection{Modeling biosystems at atomistic resolution}

\subsection{Modeling of mesoscopic and macroscopic biological systems}

\subsection{Network analysis and kinetics in biosystems}


%-----------------------------------------------------------
\section{Software and computational tools in computational biology}
%-----------------------------------------------------------
\begin{frame}{Program} 
\tableofcontents [sectionstyle=show/shaded,subsectionstyle=show/show/hide] 
\end{frame}

\subsection{General programming platforms optimal for biological problems}

\subsection{Image processing software}

\subsection{Choosing a proper software for a particular computational biological problem}

\subsection{Machine learning}

\begin{frame}{slides}

\end{frame}
%-----------------------------------------------------------
%\section*{Slidehola}
%-----------------------------------------------------------

% For those slides that might be interesting or covering advanced topics but still
%  ouside the real scope of the course.
\appendix
%-----------------------------------------------------------
\begin{frame}{Supporting slides}

\end{frame}

\end{document}

%%%%%%%%%%%%%%%%%%%%%%%%%%%%%%%%%%%%%%%%%%%%%
%%% TEMPLATES SLIDES (add more if needed) %%%
%%%%%%%%%%%%%%%%%%%%%%%%%%%%%%%%%%%%%%%%%%%%%

\begin{frame}{Title}
  \begin{columns}[c]
  %Left column  
  \column{.45\textwidth} % Left column and width
    \begin{block}{ BLOCK1}
      \begin{itemize}
        \item bla bla. 
        \item blc blc
      \end{itemize}
    \end{block}
    %Right column  
    \column{.45\textwidth} % Right column and width
    \includegraphics[width=1.0\linewidth]{fig_whatever.png}
  
  \end{columns}
\end{frame}
\note{
  \begin{itemize}
    \item note a
    \item note b
  \end{itemize}
}


